\documentclass{oci}
\usepackage[utf8]{inputenc}
\usepackage{tikz}

\definecolor{blue}{RGB}{30, 136, 229}
\definecolor{yellow}{RGB}{255, 193, 7}
\definecolor{red}{RGB}{216, 27, 96}
\definecolor{lightgray}{rgb}{0.83, 0.83, 0.83}

\newcommand{\swordred}{\hspace{0.1em}\raisebox{-0.2em}{\includegraphics[scale=0.5]{sword-red}}}
\newcommand{\swordblue}{\raisebox{-0.2em}{\includegraphics[scale=0.5]{sword-blue}}}

\title{La torre de la maldad}

\begin{document}
\begin{problemDescription}
  Como prueba final de su valentía, Valeria ha sido encomendada con la misión de infiltrarse en la
  torre de la maldad y poner fin al reinado de la liga del mal.

  La torre de la maldad tiene $N$ niveles numerados de 1 a $N$.
  En cada nivel hay un enemigo con poder $P_i$.
  Para que la misión sea satisfactoria, Valeria debe derrotar a todos los enemigos.
  El poder inicial de Valeria es $V$.
  Valeria solo puede derrotar a un enemigo si su poder $V$ es mayor o igual que el poder $P_i$
  del enemigo.
  Después de derrotar a un enemigo, este es eliminado de la torre y el nivel queda \emph{vacío}.
  Si en cualquier momento se enfrenta a un enemigo con un poder mayor que el de ella, Valeria perderá
  la misión.
  Afortunadamente, gracias a la gema de la experiencia, cada vez que Valeria derrota a un enemigo
  su poder aumenta en 1.

  Al comenzar la misión, Valeria puede infiltrarse en cualquier nivel de la torre y
  enfrentarse con el enemigo en aquel nivel.
  Una vez adentro, solo puedo moverse al siguiente nivel no vacío ya sea hacia arriba o hacia abajo.
  Nuestra heroína está interesada en saber en qué orden visitar los niveles y enfrentarse a los enemigos
  para que su misión sea exitosa.

  La siguiente imagen muestra un ejemplo para una torre con $N=5$ niveles.
  El poder inicial de Valeria es $12$.
  Por lo tanto puede decidir inicialmente infiltrarse en el nivel 3 y derrotar al enemigo cuyo poder también es 12.
  Después de derrotarlo, el poder de Valeria es 13.
  El siguiente enemigo a enfrentar debe ser el del nivel 2 el del nivel 4.
  El enemigo en el nivel 2 tiene poder 14 por lo tanto Valeria no puede derrotarlo.
  Por lo tanto está forzada a enfrentarse al enemigo en nivel 4 y cuyo poder es 13.
  El poder de Valeria es ahora 14 y las posibilidades son moverse al nivel 2 o al 5.
  Similarmente, la única alternativa es moverse al nivel 2 y enfrentar al enemigo con poder 14.
  En este punto, Valeria puede decidir derrotar al enemigo en el nivel 1 o el 5.
  Para este caso elige primero derrotar al enemigo en el nivel 1 y luego en al enemigo en el nivel 5.
  La secuencia en la que visito los niveles es la siguiente 3, 4, 2, 1, 5.
  Notar que si inicialmente Valeria hubiese decidido infiltrarse en el nivel 1, hubiese
  fallado la misión.

  \begin{center}
  \scalebox{0.8}{
    \begin{tikzpicture}
      \def\N{5}

      \begin{scope}
        \draw (0, 0) grid (1,\N);
        \foreach \y in {1, ..., \N} {
          \node[scale=0.8] at (-0.3, \y - 0.5) {\y};
        }

        \node[red, scale=0.9] at (0.5, 4.5) {\bf 15\swordred};
        \node[red, scale=0.9] at (0.5, 3.5) {\bf 13\swordred};
        \node[red, scale=0.9] at (0.5, 2.5) {\bf 12\swordred};
        \node[red, scale=0.9] at (0.5, 1.5) {\bf 14\swordred};
        \node[red, scale=0.9] at (0.5, 0.5) {\bf 11\swordred};

        \node[blue, scale=0.9] at (1.5, 2.5) {\bf \swordblue12};
      \end{scope}

      \begin{scope}[shift={(4cm, 0)}]
        \draw (0, 0) grid (1,\N);
        \foreach \y in {1, ..., \N} {
          \node[scale=0.8] at (-0.3, \y - 0.5) {\y};
        }

        \node[red, scale=0.9] at (0.5, 4.5) {\bf 15\swordred};
        \node[red, scale=0.9] at (0.5, 3.5) {\bf 13\swordred};
        \node[red, scale=0.9] at (0.5, 2.5) {};
        \node[red, scale=0.9] at (0.5, 1.5) {\bf 14\swordred};
        \node[red, scale=0.9] at (0.5, 0.5) {\bf 11\swordred};

        \node[blue, scale=0.9] at (1.5, 3.5) {\bf \swordblue13};
      \end{scope}

      \begin{scope}[shift={(8cm, 0)}]
        \draw (0, 0) grid (1,\N);
        \foreach \y in {1, ..., \N} {
          \node[scale=0.8] at (-0.3, \y - 0.5) {\y};
        }

        \node[red, scale=0.9] at (0.5, 4.5) {\bf 15\swordred};
        \node[red, scale=0.9] at (0.5, 3.5) {};
        \node[red, scale=0.9] at (0.5, 2.5) {};
        \node[red, scale=0.9] at (0.5, 1.5) {\bf 14\swordred};
        \node[red, scale=0.9] at (0.5, 0.5) {\bf 11\swordred};

        \node[blue, scale=0.9] at (1.5, 1.5) {\bf \swordblue14};
      \end{scope}

      \begin{scope}[shift={(12cm, 0)}]
        \draw (0, 0) grid (1,\N);
        \foreach \y in {1, ..., \N} {
          \node[scale=0.8] at (-0.3, \y - 0.5) {\y};
        }

        \node[red, scale=0.9] at (0.5, 4.5) {\bf 15\swordred};
        \node[red, scale=0.9] at (0.5, 3.5) {};
        \node[red, scale=0.9] at (0.5, 2.5) {};
        \node[red, scale=0.9] at (0.5, 1.5) {};
        \node[red, scale=0.9] at (0.5, 0.5) {\bf 11\swordred};

        \node[blue, scale=0.9] at (1.5, 0.5) {\bf \swordblue15};
      \end{scope}

      \begin{scope}[shift={(16cm, 0)}]
        \draw (0, 0) grid (1,\N);
        \foreach \y in {1, ..., \N} {
          \node[scale=0.8] at (-0.3, \y - 0.5) {\y};
        }

        \node[red, scale=0.9] at (0.5, 4.5) {\bf 15\swordred};
        \node[red, scale=0.9] at (0.5, 3.5) {};
        \node[red, scale=0.9] at (0.5, 2.5) {};
        \node[red, scale=0.9] at (0.5, 1.5) {};
        \node[red, scale=0.9] at (0.5, 0.5) {};

        \node[blue, scale=0.9] at (1.5, 4.5) {\bf \swordblue6};
      \end{scope}

    \end{tikzpicture}
  }
  \end{center}
\end{problemDescription}

\begin{inputDescription}
  La primera línea de la entrada contiene dos enteros $N$ y $V$ ($0 < N \leq 500$ y $0 < V \leq 10^9$)
  correspondientes respectivamente a la cantidad de niveles en la torre y el poder inicial de Valeria.
  A continuación siguen $N$ líneas cada una conteniendo un entero mayor que cero.
  El entero en la $i$-ésimo línea corresponde al poder $P_i$ del enemigo en el nivel $i$.
\end{inputDescription}

\begin{outputDescription}
  La salida debe contener línea con $N$ enteros entre 1 y $N$ correspondientes a un posible orden
  en que Valeria puede visitar los niveles para derrotar a todos los enemigos.
  En caso de haber más de un posible orden cualquiera de ellos es una respuesta válida.
  Si no existe un orden en el que sea posible derrotar a todos los enemigos, la salida debe contener
  una línea con la palabra \texttt{FALLIDA}.
\end{outputDescription}

\begin{scoreDescription}
  \subtask{100}
  Se probarán varios casos sin restricciones adicionales.
\end{scoreDescription}

\begin{sampleDescription}
\sampleIO{sample-1}
\sampleIO{sample-2}
\end{sampleDescription}

\end{document}
